%----------------------------------------------------------------------------------------
%    PACKAGES AND THEMES
%----------------------------------------------------------------------------------------
\documentclass[aspectratio=169,xcolor=dvipsnames]{beamer}
\makeatletter
\def\input@path{{../BeamerTemplate/theme/}}
\makeatother
\usetheme{CleanEasy}
\usepackage[utf8]{inputenc}
\usepackage{lmodern}
\usepackage[T1]{fontenc}
\usepackage{fix-cm}
\usepackage{amsmath}
\usepackage{mathtools}
\usepackage{listings}
\usepackage{xcolor}
\usepackage{hyperref}
\usepackage{graphicx}
\usepackage{booktabs}
\usepackage{tikz}
\usetikzlibrary{positioning, shapes, arrows, calc, decorations.pathreplacing, arrows.meta, backgrounds, patterns, overlay-beamer-styles}
\usepackage{etoolbox}
\usepackage{animate}

%----------------------------------------------------------------------------------------
%    LAYOUT CONFIGURATION
%----------------------------------------------------------------------------------------
\input{../BeamerTemplate/configs/configs}

%----------------------------------------------------------------------------------------
%    TITLE PAGE CONFIGURATION
%----------------------------------------------------------------------------------------
\title{Algorithms with Data Structures}
\subtitle{Sorting, Searching, Graph Algorithms, and Optimization Patterns}
\author{Minseok Jeon}
\institute{DGIST}
\date{\today}

\begin{document}

\begin{frame}
  \titlepage
\end{frame}

\begin{frame}{Outline}
  \tableofcontents
\end{frame}

% Introduction
\section{Introduction}

\begin{frame}{Course Overview}
  \begin{itemize}
    \item \textbf{Sorting Algorithms}: Comparison-based and non-comparison-based
    \item \textbf{Searching Algorithms}: Linear, binary, and advanced techniques
    \item \textbf{Graph Algorithms}: Shortest path and minimum spanning tree
    \item \textbf{Dynamic Programming}: Optimization with data structures
    \item \textbf{Complexity-Driven Design}: Choosing the right approach
    \item \textbf{Practical Patterns}: Two pointers, sliding window, and more
  \end{itemize}
\end{frame}

\begin{frame}{Why Algorithms + Data Structures?}
  \begin{block}{Synergy}
    Algorithms are meaningless without efficient data structures
  \end{block}

  \begin{itemize}
    \item Data structures enable efficient algorithm implementation
    \item Algorithm choice depends on data structure characteristics
    \item Complexity analysis requires understanding both
    \item Real-world performance depends on the combination
  \end{itemize}
\end{frame}

%% Sorting Algorithms
\section{Sorting Algorithms}

\begin{frame}{Sorting: The Foundation}
  \begin{block}{What is Sorting?}
    Arranging elements in a specific order (ascending/descending)
  \end{block}

  \textbf{Two Categories:}
  \begin{itemize}
    \item \textbf{Comparison-based}: Compare elements to determine order
    \item \textbf{Non-comparison-based}: Use element properties (e.g., digits)
  \end{itemize}

  \vspace{0.3cm}
  \textbf{Why It Matters:}
  \begin{itemize}
    \item Enables binary search ($O(\log n)$ vs $O(n)$)
    \item Foundation for many algorithms
    \item Common interview topic
  \end{itemize}
\end{frame}

\subsection{Comparison-Based Sorting}

\begin{frame}{Bubble Sort}
  \begin{columns}
    \begin{column}{0.5\textwidth}
      \textbf{Algorithm:}
      \begin{enumerate}
        \item Compare adjacent elements
        \item Swap if out of order
        \item Repeat until sorted
      \end{enumerate}

      \vspace{0.3cm}
      \textbf{Characteristics:}
      \begin{itemize}
        \item Time: $O(n^2)$
        \item Space: $O(1)$
        \item Stable: Yes
      \end{itemize}
    \end{column}

    \begin{column}{0.5\textwidth}
      \textbf{Example:}
      \begin{align*}
        &[5, 2, 8, 1, 9] \\
        &[2, 5, 1, 8, 9] \\
        &[2, 1, 5, 8, 9] \\
        &[1, 2, 5, 8, 9]
      \end{align*}

      \vspace{0.3cm}
      \textbf{Use Case:}
      \begin{itemize}
        \item Small datasets
        \item Educational purposes
        \item Nearly sorted data
      \end{itemize}
    \end{column}
  \end{columns}
\end{frame}

\begin{frame}[fragile]{Bubble Sort: Implementation}
\begin{lstlisting}[language=Python]
def bubble_sort(arr):
    n = len(arr)
    for i in range(n):
        swapped = False
        for j in range(0, n - i - 1):
            if arr[j] > arr[j + 1]:
                arr[j], arr[j + 1] = arr[j + 1], arr[j]
                swapped = True
        if not swapped:
            break  # Early termination
    return arr
\end{lstlisting}
\end{frame}

\begin{frame}{Selection Sort}
  \begin{columns}
    \begin{column}{0.5\textwidth}
      \textbf{Algorithm:}
      \begin{enumerate}
        \item Find minimum element
        \item Swap with first unsorted position
        \item Repeat for remaining elements
      \end{enumerate}

      \vspace{0.3cm}
      \textbf{Characteristics:}
      \begin{itemize}
        \item Time: $O(n^2)$
        \item Space: $O(1)$
        \item Stable: No
      \end{itemize}
    \end{column}

    \begin{column}{0.5\textwidth}
      \textbf{Example:}
      \begin{align*}
        &[5, 2, 8, 1, 9] \\
        &[1, 2, 8, 5, 9] \\
        &[1, 2, 8, 5, 9] \\
        &[1, 2, 5, 8, 9]
      \end{align*}

      \vspace{0.3cm}
      \textbf{Use Case:}
      \begin{itemize}
        \item Minimal swaps needed
        \item Small datasets
      \end{itemize}
    \end{column}
  \end{columns}
\end{frame}

\begin{frame}{Insertion Sort}
  \begin{columns}
    \begin{column}{0.5\textwidth}
      \textbf{Algorithm:}
      \begin{enumerate}
        \item Take next element
        \item Insert into sorted portion
        \item Shift elements as needed
      \end{enumerate}

      \vspace{0.3cm}
      \textbf{Characteristics:}
      \begin{itemize}
        \item Time: $O(n^2)$ average, $O(n)$ best
        \item Space: $O(1)$
        \item Stable: Yes
        \item Adaptive: Yes
      \end{itemize}
    \end{column}

    \begin{column}{0.5\textwidth}
      \textbf{Advantages:}
      \begin{itemize}
        \item Efficient for small datasets
        \item Excellent for nearly sorted data
        \item Online algorithm
        \item Used in Timsort
      \end{itemize}

      \vspace{0.3cm}
      \textbf{Use Case:}
      \begin{itemize}
        \item Small arrays ($n < 50$)
        \item Nearly sorted data
        \item Streaming data
      \end{itemize}
    \end{column}
  \end{columns}
\end{frame}

\begin{frame}{Merge Sort}
  \begin{columns}
    \begin{column}{0.5\textwidth}
      \textbf{Algorithm (Divide \& Conquer):}
      \begin{enumerate}
        \item Divide array in half
        \item Recursively sort each half
        \item Merge sorted halves
      \end{enumerate}

      \vspace{0.3cm}
      \textbf{Characteristics:}
      \begin{itemize}
        \item Time: $O(n \log n)$ always
        \item Space: $O(n)$
        \item Stable: Yes
      \end{itemize}
    \end{column}

    \begin{column}{0.5\textwidth}
      \textbf{Advantages:}
      \begin{itemize}
        \item Guaranteed $O(n \log n)$
        \item Stable sorting
        \item Parallelizable
        \item Good for linked lists
      \end{itemize}

      \vspace{0.3cm}
      \textbf{Use Case:}
      \begin{itemize}
        \item External sorting
        \item Linked lists
        \item Stability required
      \end{itemize}
    \end{column}
  \end{columns}
\end{frame}

\begin{frame}{Quick Sort}
  \begin{columns}
    \begin{column}{0.5\textwidth}
      \textbf{Algorithm (Divide \& Conquer):}
      \begin{enumerate}
        \item Choose pivot element
        \item Partition around pivot
        \item Recursively sort partitions
      \end{enumerate}

      \vspace{0.3cm}
      \textbf{Characteristics:}
      \begin{itemize}
        \item Time: $O(n \log n)$ avg, $O(n^2)$ worst
        \item Space: $O(\log n)$ stack
        \item Stable: No
        \item In-place: Yes
      \end{itemize}
    \end{column}

    \begin{column}{0.5\textwidth}
      \textbf{Pivot Selection:}
      \begin{itemize}
        \item First/Last: Simple but risky
        \item Random: Better average case
        \item Median-of-three: Balanced
      \end{itemize}

      \vspace{0.3cm}
      \textbf{Use Case:}
      \begin{itemize}
        \item General-purpose (most libraries)
        \item Cache-efficient
        \item In-place sorting needed
      \end{itemize}
    \end{column}
  \end{columns}
\end{frame}

\begin{frame}{Heap Sort}
  \begin{columns}
    \begin{column}{0.5\textwidth}
      \textbf{Algorithm:}
      \begin{enumerate}
        \item Build max heap
        \item Swap root with last element
        \item Heapify remaining elements
        \item Repeat
      \end{enumerate}

      \vspace{0.3cm}
      \textbf{Characteristics:}
      \begin{itemize}
        \item Time: $O(n \log n)$ always
        \item Space: $O(1)$
        \item Stable: No
        \item In-place: Yes
      \end{itemize}
    \end{column}

    \begin{column}{0.5\textwidth}
      \textbf{Advantages:}
      \begin{itemize}
        \item Guaranteed $O(n \log n)$
        \item In-place sorting
        \item No worst-case degradation
      \end{itemize}

      \vspace{0.3cm}
      \textbf{Use Case:}
      \begin{itemize}
        \item Memory-constrained systems
        \item Real-time systems
        \item Embedded systems
      \end{itemize}
    \end{column}
  \end{columns}
\end{frame}

\begin{frame}{Comparison-Based Sorting: Summary}
  \begin{table}
    \centering
    \small
    \begin{tabular}{lccccc}
      \hline
      \textbf{Algorithm} & \textbf{Best} & \textbf{Average} & \textbf{Worst} & \textbf{Space} & \textbf{Stable} \\
      \hline
      Bubble & $O(n)$ & $O(n^2)$ & $O(n^2)$ & $O(1)$ & Yes \\
      Selection & $O(n^2)$ & $O(n^2)$ & $O(n^2)$ & $O(1)$ & No \\
      Insertion & $O(n)$ & $O(n^2)$ & $O(n^2)$ & $O(1)$ & Yes \\
      Merge & $O(n \log n)$ & $O(n \log n)$ & $O(n \log n)$ & $O(n)$ & Yes \\
      Quick & $O(n \log n)$ & $O(n \log n)$ & $O(n^2)$ & $O(\log n)$ & No \\
      Heap & $O(n \log n)$ & $O(n \log n)$ & $O(n \log n)$ & $O(1)$ & No \\
      \hline
    \end{tabular}
  \end{table}

  \vspace{0.3cm}
  \textbf{Lower Bound}: $\Omega(n \log n)$ for comparison-based sorting
\end{frame}

\subsection{Non-Comparison-Based Sorting}

\begin{frame}{Counting Sort}
  \begin{columns}
    \begin{column}{0.5\textwidth}
      \textbf{Algorithm:}
      \begin{enumerate}
        \item Count occurrences of each value
        \item Calculate cumulative counts
        \item Place elements in sorted order
      \end{enumerate}

      \vspace{0.3cm}
      \textbf{Characteristics:}
      \begin{itemize}
        \item Time: $O(n + k)$
        \item Space: $O(k)$
        \item Stable: Yes
        \item $k$ = range of input
      \end{itemize}
    \end{column}

    \begin{column}{0.5\textwidth}
      \textbf{Constraints:}
      \begin{itemize}
        \item Integer keys only
        \item Known range required
        \item Range must be reasonable
      \end{itemize}

      \vspace{0.3cm}
      \textbf{Use Case:}
      \begin{itemize}
        \item Small integer range
        \item Ages, grades, scores
        \item Subroutine for radix sort
      \end{itemize}
    \end{column}
  \end{columns}
\end{frame}

\begin{frame}{Radix Sort}
  \begin{columns}
    \begin{column}{0.5\textwidth}
      \textbf{Algorithm:}
      \begin{enumerate}
        \item Sort by least significant digit
        \item Move to next digit
        \item Repeat until all digits processed
      \end{enumerate}

      \vspace{0.3cm}
      \textbf{Characteristics:}
      \begin{itemize}
        \item Time: $O(d \cdot (n + k))$
        \item Space: $O(n + k)$
        \item Stable: Yes
        \item $d$ = number of digits
      \end{itemize}
    \end{column}

    \begin{column}{0.5\textwidth}
      \textbf{Variants:}
      \begin{itemize}
        \item LSD: Least significant digit first
        \item MSD: Most significant digit first
      \end{itemize}

      \vspace{0.3cm}
      \textbf{Use Case:}
      \begin{itemize}
        \item Fixed-length integers
        \item Strings of equal length
        \item Large datasets with small digit count
      \end{itemize}
    \end{column}
  \end{columns}
\end{frame}

\begin{frame}{Bucket Sort}
  \begin{columns}
    \begin{column}{0.5\textwidth}
      \textbf{Algorithm:}
      \begin{enumerate}
        \item Distribute elements into buckets
        \item Sort each bucket individually
        \item Concatenate sorted buckets
      \end{enumerate}

      \vspace{0.3cm}
      \textbf{Characteristics:}
      \begin{itemize}
        \item Time: $O(n + k)$ average
        \item Space: $O(n + k)$
        \item Stable: Depends on sub-sort
      \end{itemize}
    \end{column}

    \begin{column}{0.5\textwidth}
      \textbf{Requirements:}
      \begin{itemize}
        \item Uniform distribution
        \item Known input range
      \end{itemize}

      \vspace{0.3cm}
      \textbf{Use Case:}
      \begin{itemize}
        \item Uniformly distributed data
        \item Floating-point numbers [0, 1)
        \item External sorting
      \end{itemize}
    \end{column}
  \end{columns}
\end{frame}

\begin{frame}{Non-Comparison Sorting: Summary}
  \begin{table}
    \centering
    \small
    \begin{tabular}{lcccc}
      \hline
      \textbf{Algorithm} & \textbf{Time} & \textbf{Space} & \textbf{Stable} & \textbf{Constraints} \\
      \hline
      Counting & $O(n + k)$ & $O(k)$ & Yes & Integer, known range \\
      Radix & $O(d(n + k))$ & $O(n + k)$ & Yes & Fixed-length keys \\
      Bucket & $O(n + k)$ & $O(n + k)$ & Varies & Uniform distribution \\
      \hline
    \end{tabular}
  \end{table}

  \vspace{0.3cm}
  \begin{block}{Key Insight}
    Non-comparison sorts can beat the $\Omega(n \log n)$ lower bound by exploiting specific properties of the input data
  \end{block}
\end{frame}

%% Searching Algorithms
\section{Searching Algorithms}

\begin{frame}{Searching: Finding Elements Efficiently}
  \begin{block}{Goal}
    Find the position or existence of a target element in a dataset
  \end{block}

  \vspace{0.3cm}
  \textbf{Categories:}
  \begin{itemize}
    \item \textbf{Basic}: Linear search
    \item \textbf{Sorted Array}: Binary search and variants
    \item \textbf{Hash-based}: Constant time lookup
  \end{itemize}

  \vspace{0.3cm}
  \textbf{Trade-offs:}
  \begin{itemize}
    \item Preprocessing vs. query time
    \item Space vs. time complexity
    \item Data structure requirements
  \end{itemize}
\end{frame}

\subsection{Basic Search}

\begin{frame}{Linear Search}
  \begin{columns}
    \begin{column}{0.5\textwidth}
      \textbf{Algorithm:}
      \begin{enumerate}
        \item Check each element sequentially
        \item Return index if found
        \item Return -1 if not found
      \end{enumerate}

      \vspace{0.3cm}
      \textbf{Characteristics:}
      \begin{itemize}
        \item Time: $O(n)$
        \item Space: $O(1)$
        \item No preprocessing needed
      \end{itemize}
    \end{column}

    \begin{column}{0.5\textwidth}
      \textbf{Advantages:}
      \begin{itemize}
        \item Works on unsorted data
        \item Simple implementation
        \item No extra space
      \end{itemize}

      \vspace{0.3cm}
      \textbf{Use Case:}
      \begin{itemize}
        \item Small datasets
        \item Unsorted data
        \item One-time searches
      \end{itemize}
    \end{column}
  \end{columns}
\end{frame}

\subsection{Sorted Array Search}

\begin{frame}{Binary Search}
  \begin{columns}
    \begin{column}{0.5\textwidth}
      \textbf{Algorithm (Divide \& Conquer):}
      \begin{enumerate}
        \item Compare target with middle element
        \item Eliminate half of search space
        \item Repeat until found or exhausted
      \end{enumerate}

      \vspace{0.3cm}
      \textbf{Characteristics:}
      \begin{itemize}
        \item Time: $O(\log n)$
        \item Space: $O(1)$ iterative, $O(\log n)$ recursive
        \item \textbf{Requires sorted array}
      \end{itemize}
    \end{column}

    \begin{column}{0.5\textwidth}
      \textbf{Variants:}
      \begin{itemize}
        \item Find first occurrence
        \item Find last occurrence
        \item Count occurrences
        \item Find insertion position
      \end{itemize}

      \vspace{0.3cm}
      \textbf{Use Case:}
      \begin{itemize}
        \item Sorted datasets
        \item Repeated queries
        \item Most common search algorithm
      \end{itemize}
    \end{column}
  \end{columns}
\end{frame}

\begin{frame}[fragile]{Binary Search: Template}
\begin{lstlisting}[language=Python]
def binary_search(arr, target):
    left, right = 0, len(arr) - 1

    while left <= right:
        mid = left + (right - left) // 2

        if arr[mid] == target:
            return mid
        elif arr[mid] < target:
            left = mid + 1
        else:
            right = mid - 1

    return -1  # Not found
\end{lstlisting}
\end{frame}

\begin{frame}{Interpolation Search}
  \begin{columns}
    \begin{column}{0.5\textwidth}
      \textbf{Algorithm:}
      \begin{enumerate}
        \item Estimate position using interpolation
        \item Check estimated position
        \item Adjust search range
      \end{enumerate}

      \vspace{0.3cm}
      \textbf{Position Formula:}
      \begin{align*}
        \text{pos} = \text{low} + \frac{(x - arr[\text{low}])}{(arr[\text{high}] - arr[\text{low}])} \times (\text{high} - \text{low})
      \end{align*}
    \end{column}

    \begin{column}{0.5\textwidth}
      \textbf{Characteristics:}
      \begin{itemize}
        \item Time: $O(\log \log n)$ avg, $O(n)$ worst
        \item Space: $O(1)$
        \item Requires uniform distribution
      \end{itemize}

      \vspace{0.3cm}
      \textbf{Use Case:}
      \begin{itemize}
        \item Uniformly distributed data
        \item Large sorted datasets
        \item Phone books, dictionaries
      \end{itemize}
    \end{column}
  \end{columns}
\end{frame}

\begin{frame}{Exponential Search}
  \begin{columns}
    \begin{column}{0.5\textwidth}
      \textbf{Algorithm:}
      \begin{enumerate}
        \item Find range with exponential growth
        \item Perform binary search in range
      \end{enumerate}

      \vspace{0.3cm}
      \textbf{Characteristics:}
      \begin{itemize}
        \item Time: $O(\log n)$
        \item Space: $O(1)$
        \item Better for unbounded arrays
      \end{itemize}
    \end{column}

    \begin{column}{0.5\textwidth}
      \textbf{Advantages:}
      \begin{itemize}
        \item Works on unbounded/infinite arrays
        \item Better than binary when target is near start
      \end{itemize}

      \vspace{0.3cm}
      \textbf{Use Case:}
      \begin{itemize}
        \item Unbounded search space
        \item Target likely near beginning
        \item Infinite streams
      \end{itemize}
    \end{column}
  \end{columns}
\end{frame}

\begin{frame}{Jump Search}
  \begin{columns}
    \begin{column}{0.5\textwidth}
      \textbf{Algorithm:}
      \begin{enumerate}
        \item Jump by fixed block size $\sqrt{n}$
        \item Find block containing target
        \item Linear search within block
      \end{enumerate}

      \vspace{0.3cm}
      \textbf{Characteristics:}
      \begin{itemize}
        \item Time: $O(\sqrt{n})$
        \item Space: $O(1)$
        \item Optimal jump: $\sqrt{n}$
      \end{itemize}
    \end{column}

    \begin{column}{0.5\textwidth}
      \textbf{Advantages:}
      \begin{itemize}
        \item Better than linear for sorted arrays
        \item Fewer comparisons than binary
        \item Good for jumping in physical storage
      \end{itemize}

      \vspace{0.3cm}
      \textbf{Use Case:}
      \begin{itemize}
        \item Systems where jumping back is costly
        \item Disk-based systems
      \end{itemize}
    \end{column}
  \end{columns}
\end{frame}

\begin{frame}{Ternary Search}
  \begin{columns}
    \begin{column}{0.5\textwidth}
      \textbf{Algorithm:}
      \begin{enumerate}
        \item Divide range into three parts
        \item Determine which third contains target
        \item Recursively search that third
      \end{enumerate}

      \vspace{0.3cm}
      \textbf{Characteristics:}
      \begin{itemize}
        \item Time: $O(\log_3 n)$
        \item More comparisons per iteration than binary
      \end{itemize}
    \end{column}

    \begin{column}{0.5\textwidth}
      \textbf{Comparison with Binary:}
      \begin{itemize}
        \item $\log_3 n < \log_2 n$ (fewer iterations)
        \item But 2 comparisons per iteration vs 1
        \item Binary is generally faster
      \end{itemize}

      \vspace{0.3cm}
      \textbf{Use Case:}
      \begin{itemize}
        \item Finding max/min of unimodal function
        \item Optimization problems
      \end{itemize}
    \end{column}
  \end{columns}
\end{frame}

\subsection{Hash-Based Search}

\begin{frame}{Hash Table Search}
  \begin{columns}
    \begin{column}{0.5\textwidth}
      \textbf{Approach:}
      \begin{enumerate}
        \item Hash key to index
        \item Access bucket at index
        \item Handle collisions if needed
      \end{enumerate}

      \vspace{0.3cm}
      \textbf{Characteristics:}
      \begin{itemize}
        \item Time: $O(1)$ avg, $O(n)$ worst
        \item Space: $O(n)$
        \item Requires hash function
      \end{itemize}
    \end{column}

    \begin{column}{0.5\textwidth}
      \textbf{Trade-offs:}
      \begin{itemize}
        \item Space for time
        \item No ordering maintained
        \item Hash collisions possible
      \end{itemize}

      \vspace{0.3cm}
      \textbf{Use Case:}
      \begin{itemize}
        \item Frequent lookups
        \item Unordered data
        \item Database indexing
        \item Caching
      \end{itemize}
    \end{column}
  \end{columns}
\end{frame}

\begin{frame}{Searching Algorithms: Summary}
  \begin{table}
    \centering
    \small
    \begin{tabular}{lccl}
      \hline
      \textbf{Algorithm} & \textbf{Time} & \textbf{Space} & \textbf{Requirement} \\
      \hline
      Linear & $O(n)$ & $O(1)$ & None \\
      Binary & $O(\log n)$ & $O(1)$ & Sorted \\
      Interpolation & $O(\log \log n)$ & $O(1)$ & Sorted + Uniform \\
      Exponential & $O(\log n)$ & $O(1)$ & Sorted + Unbounded \\
      Jump & $O(\sqrt{n})$ & $O(1)$ & Sorted \\
      Ternary & $O(\log_3 n)$ & $O(1)$ & Sorted / Unimodal \\
      Hash Table & $O(1)$ avg & $O(n)$ & Hash function \\
      \hline
    \end{tabular}
  \end{table}
\end{frame}

%% Graph Algorithms
\section{Graph Algorithms}

\begin{frame}{Graph Algorithms Overview}
  \begin{block}{Graph Problems}
    Graphs model relationships between entities
  \end{block}

  \vspace{0.3cm}
  \textbf{Core Categories:}
  \begin{itemize}
    \item \textbf{Shortest Path}: Find minimum-cost path between vertices
    \item \textbf{Minimum Spanning Tree (MST)}: Connect all vertices with minimum total edge weight
  \end{itemize}

  \vspace{0.3cm}
  \textbf{Applications:}
  \begin{itemize}
    \item Network routing, GPS navigation
    \item Social networks, recommendation systems
    \item Circuit design, network infrastructure
  \end{itemize}
\end{frame}

\subsection{Shortest Path Algorithms}

\begin{frame}{Dijkstra's Algorithm}
  \begin{columns}
    \begin{column}{0.5\textwidth}
      \textbf{Algorithm (Greedy):}
      \begin{enumerate}
        \item Initialize distances (source = 0, others = $\infty$)
        \item Extract minimum distance vertex
        \item Relax edges from that vertex
        \item Repeat until all processed
      \end{enumerate}

      \vspace{0.3cm}
      \textbf{Data Structure:}
      \begin{itemize}
        \item Priority queue (min-heap)
      \end{itemize}
    \end{column}

    \begin{column}{0.5\textwidth}
      \textbf{Characteristics:}
      \begin{itemize}
        \item Time: $O((V + E) \log V)$ with heap
        \item Space: $O(V)$
        \item \textbf{Non-negative weights only}
        \item Single-source shortest path
      \end{itemize}

      \vspace{0.3cm}
      \textbf{Use Case:}
      \begin{itemize}
        \item GPS navigation
        \item Network routing (OSPF)
        \item Most common shortest path
      \end{itemize}
    \end{column}
  \end{columns}
\end{frame}

\begin{frame}[fragile]{Dijkstra's Algorithm: Implementation}
\begin{lstlisting}[language=Python]
import heapq

def dijkstra(graph, start):
    distances = {v: float('inf') for v in graph}
    distances[start] = 0
    pq = [(0, start)]  # (distance, vertex)

    while pq:
        curr_dist, u = heapq.heappop(pq)
        if curr_dist > distances[u]:
            continue

        for v, weight in graph[u]:
            distance = curr_dist + weight
            if distance < distances[v]:
                distances[v] = distance
                heapq.heappush(pq, (distance, v))

    return distances
\end{lstlisting}
\end{frame}

\begin{frame}{Bellman-Ford Algorithm}
  \begin{columns}
    \begin{column}{0.5\textwidth}
      \textbf{Algorithm (Dynamic Programming):}
      \begin{enumerate}
        \item Initialize distances
        \item Relax all edges $V-1$ times
        \item Check for negative cycles
      \end{enumerate}

      \vspace{0.3cm}
      \textbf{Characteristics:}
      \begin{itemize}
        \item Time: $O(VE)$
        \item Space: $O(V)$
        \item Handles negative weights
        \item Detects negative cycles
      \end{itemize}
    \end{column}

    \begin{column}{0.5\textwidth}
      \textbf{Advantages over Dijkstra:}
      \begin{itemize}
        \item Negative edge weights OK
        \item Detects negative cycles
        \item Simpler implementation
      \end{itemize}

      \vspace{0.3cm}
      \textbf{Use Case:}
      \begin{itemize}
        \item Graphs with negative weights
        \item Detecting arbitrage opportunities
        \item Distance vector routing
      \end{itemize}
    \end{column}
  \end{columns}
\end{frame}

\begin{frame}{Floyd-Warshall Algorithm}
  \begin{columns}
    \begin{column}{0.5\textwidth}
      \textbf{Algorithm (Dynamic Programming):}
      \begin{enumerate}
        \item Initialize distance matrix
        \item For each intermediate vertex $k$
        \item Try path through $k$
        \item Update if shorter
      \end{enumerate}

      \vspace{0.3cm}
      \textbf{DP Formula:}
      \begin{align*}
        \text{dist}[i][j] = \min(&\text{dist}[i][j], \\
        &\text{dist}[i][k] + \text{dist}[k][j])
      \end{align*}
    \end{column}

    \begin{column}{0.5\textwidth}
      \textbf{Characteristics:}
      \begin{itemize}
        \item Time: $O(V^3)$
        \item Space: $O(V^2)$
        \item All-pairs shortest paths
        \item Handles negative weights
      \end{itemize}

      \vspace{0.3cm}
      \textbf{Use Case:}
      \begin{itemize}
        \item Dense graphs
        \item All-pairs distances needed
        \item Transitive closure
      \end{itemize}
    \end{column}
  \end{columns}
\end{frame}

\begin{frame}{A* Search Algorithm}
  \begin{columns}
    \begin{column}{0.5\textwidth}
      \textbf{Algorithm (Informed Search):}
      \begin{enumerate}
        \item Use heuristic function $h(n)$
        \item Evaluate $f(n) = g(n) + h(n)$
        \item Expand most promising node
      \end{enumerate}

      \vspace{0.3cm}
      \textbf{Components:}
      \begin{itemize}
        \item $g(n)$: Cost from start to $n$
        \item $h(n)$: Estimated cost from $n$ to goal
        \item $f(n)$: Total estimated cost
      \end{itemize}
    \end{column}

    \begin{column}{0.5\textwidth}
      \textbf{Heuristic Properties:}
      \begin{itemize}
        \item \textbf{Admissible}: $h(n) \le$ true cost
        \item \textbf{Consistent}: $h(n) \le c(n,n') + h(n')$
      \end{itemize}

      \vspace{0.3cm}
      \textbf{Use Case:}
      \begin{itemize}
        \item Pathfinding in games
        \item Robotics navigation
        \item GPS with traffic
      \end{itemize}
    \end{column}
  \end{columns}
\end{frame}

\begin{frame}{Shortest Path: Comparison}
  \begin{table}
    \centering
    \small
    \begin{tabular}{lcccc}
      \hline
      \textbf{Algorithm} & \textbf{Time} & \textbf{Type} & \textbf{Negative?} & \textbf{Use Case} \\
      \hline
      Dijkstra & $O((V+E) \log V)$ & Single & No & General \\
      Bellman-Ford & $O(VE)$ & Single & Yes & Negative weights \\
      Floyd-Warshall & $O(V^3)$ & All-pairs & Yes & Dense graphs \\
      A* & Varies & Single & No & Heuristic available \\
      \hline
    \end{tabular}
  \end{table}

  \vspace{0.3cm}
  \begin{block}{Selection Guide}
    \begin{itemize}
      \item Non-negative weights $\to$ Dijkstra
      \item Negative weights $\to$ Bellman-Ford
      \item All pairs needed $\to$ Floyd-Warshall
      \item Known goal + heuristic $\to$ A*
    \end{itemize}
  \end{block}
\end{frame}

\subsection{Minimum Spanning Tree (MST)}

\begin{frame}{Minimum Spanning Tree: Definition}
  \begin{block}{MST Problem}
    Find a tree that connects all vertices with minimum total edge weight
  \end{block}

  \vspace{0.3cm}
  \textbf{Properties:}
  \begin{itemize}
    \item Connects all $V$ vertices
    \item Has exactly $V-1$ edges
    \item Acyclic (no cycles)
    \item Minimum total weight
  \end{itemize}

  \vspace{0.3cm}
  \textbf{Applications:}
  \begin{itemize}
    \item Network design (minimize cable length)
    \item Cluster analysis
    \item Approximation algorithms
  \end{itemize}
\end{frame}

\begin{frame}{Kruskal's Algorithm}
  \begin{columns}
    \begin{column}{0.5\textwidth}
      \textbf{Algorithm (Greedy):}
      \begin{enumerate}
        \item Sort all edges by weight
        \item For each edge (increasing weight):
        \begin{itemize}
          \item If doesn't form cycle, add it
        \end{itemize}
        \item Stop when $V-1$ edges added
      \end{enumerate}

      \vspace{0.3cm}
      \textbf{Data Structure:}
      \begin{itemize}
        \item Union-Find (Disjoint Set)
      \end{itemize}
    \end{column}

    \begin{column}{0.5\textwidth}
      \textbf{Characteristics:}
      \begin{itemize}
        \item Time: $O(E \log E)$ or $O(E \log V)$
        \item Space: $O(V)$ for Union-Find
        \item Edge-based approach
      \end{itemize}

      \vspace{0.3cm}
      \textbf{Use Case:}
      \begin{itemize}
        \item Sparse graphs ($E \ll V^2$)
        \item Edge list representation
        \item Parallel implementation possible
      \end{itemize}
    \end{column}
  \end{columns}
\end{frame}

\begin{frame}{Prim's Algorithm}
  \begin{columns}
    \begin{column}{0.5\textwidth}
      \textbf{Algorithm (Greedy):}
      \begin{enumerate}
        \item Start with arbitrary vertex
        \item Add minimum edge to tree
        \item Expand tree vertex by vertex
        \item Stop when all vertices included
      \end{enumerate}

      \vspace{0.3cm}
      \textbf{Data Structure:}
      \begin{itemize}
        \item Priority queue (min-heap)
      \end{itemize}
    \end{column}

    \begin{column}{0.5\textwidth}
      \textbf{Characteristics:}
      \begin{itemize}
        \item Time: $O((V + E) \log V)$ with heap
        \item Space: $O(V)$
        \item Vertex-based approach
      \end{itemize}

      \vspace{0.3cm}
      \textbf{Use Case:}
      \begin{itemize}
        \item Dense graphs ($E \approx V^2$)
        \item Adjacency matrix representation
        \item Similar to Dijkstra
      \end{itemize}
    \end{column}
  \end{columns}
\end{frame}

\begin{frame}{MST: Kruskal vs Prim}
  \begin{table}
    \centering
    \begin{tabular}{lcc}
      \hline
      \textbf{Aspect} & \textbf{Kruskal} & \textbf{Prim} \\
      \hline
      Approach & Edge-based & Vertex-based \\
      Data Structure & Union-Find & Priority Queue \\
      Time Complexity & $O(E \log E)$ & $O((V+E) \log V)$ \\
      Best For & Sparse graphs & Dense graphs \\
      Graph Rep. & Edge list & Adjacency list/matrix \\
      \hline
    \end{tabular}
  \end{table}

  \vspace{0.3cm}
  \begin{block}{Selection Guide}
    \begin{itemize}
      \item Sparse graph ($E \ll V^2$) $\to$ Kruskal
      \item Dense graph ($E \approx V^2$) $\to$ Prim
      \item Both guarantee optimal MST
    \end{itemize}
  \end{block}
\end{frame}

%% Dynamic Programming
\section{Dynamic Programming with Data Structures}

\begin{frame}{Dynamic Programming Overview}
  \begin{block}{Core Idea}
    Break problem into overlapping subproblems, store results to avoid recomputation
  \end{block}

  \vspace{0.3cm}
  \textbf{Key Properties:}
  \begin{itemize}
    \item \textbf{Optimal Substructure}: Optimal solution contains optimal solutions to subproblems
    \item \textbf{Overlapping Subproblems}: Same subproblems solved multiple times
  \end{itemize}

  \vspace{0.3cm}
  \textbf{Data Structure's Role:}
  \begin{itemize}
    \item \textbf{Memoization}: Hash table for caching
    \item \textbf{Tabulation}: Arrays for bottom-up
    \item \textbf{State Optimization}: Specialized structures
  \end{itemize}
\end{frame}

\begin{frame}{DP with Arrays}
  \begin{columns}
    \begin{column}{0.5\textwidth}
      \textbf{Classic Problems:}
      \begin{itemize}
        \item Fibonacci numbers
        \item Longest Increasing Subsequence
        \item Maximum subarray sum
        \item Edit distance
      \end{itemize}

      \vspace{0.3cm}
      \textbf{Pattern:}
      \begin{enumerate}
        \item Define state: \texttt{dp[i]}
        \item Base cases
        \item Recurrence relation
        \item Compute bottom-up
      \end{enumerate}
    \end{column}

    \begin{column}{0.5\textwidth}
      \textbf{Example: Fibonacci}
      \begin{align*}
        &\text{dp}[0] = 0 \\
        &\text{dp}[1] = 1 \\
        &\text{dp}[i] = \text{dp}[i-1] + \text{dp}[i-2]
      \end{align*}

      \vspace{0.3cm}
      \textbf{Space Optimization:}
      \begin{itemize}
        \item Only need last 2 values
        \item $O(n) \to O(1)$ space
      \end{itemize}
    \end{column}
  \end{columns}
\end{frame}

\begin{frame}{DP with Hash Tables}
  \begin{columns}
    \begin{column}{0.5\textwidth}
      \textbf{Use Cases:}
      \begin{itemize}
        \item State space is sparse
        \item Multi-dimensional state
        \item State is complex (tuple, string)
      \end{itemize}

      \vspace{0.3cm}
      \textbf{Advantages:}
      \begin{itemize}
        \item Only store computed states
        \item Flexible state representation
        \item Easy memoization
      \end{itemize}
    \end{column}

    \begin{column}{0.5\textwidth}
      \textbf{Example: Word Break}
      \begin{itemize}
        \item State: remaining substring
        \item Hash table maps substring $\to$ boolean
        \item $O(n^2)$ time, $O(n)$ space
      \end{itemize}

      \vspace{0.3cm}
      \textbf{Pattern:}
      \begin{enumerate}
        \item Check if state cached
        \item If not, compute recursively
        \item Cache result before returning
      \end{enumerate}
    \end{column}
  \end{columns}
\end{frame}

\begin{frame}{DP with Trees}
  \begin{columns}
    \begin{column}{0.5\textwidth}
      \textbf{Tree DP Problems:}
      \begin{itemize}
        \item Tree diameter
        \item House robber on tree
        \item Maximum path sum
        \item Subtree queries
      \end{itemize}

      \vspace{0.3cm}
      \textbf{Pattern:}
      \begin{enumerate}
        \item DFS traversal
        \item Combine child results
        \item Return value for parent
      \end{enumerate}
    \end{column}

    \begin{column}{0.5\textwidth}
      \textbf{Example: Tree Diameter}
      \begin{itemize}
        \item State: max depth from node
        \item Combine: max of two child depths
        \item Answer: max sum of two child depths
      \end{itemize}

      \vspace{0.3cm}
      \textbf{Complexity:}
      \begin{itemize}
        \item Time: $O(n)$ (visit each node once)
        \item Space: $O(h)$ recursion stack
      \end{itemize}
    \end{column}
  \end{columns}
\end{frame}

\begin{frame}{DP with Graphs}
  \begin{columns}
    \begin{column}{0.5\textwidth}
      \textbf{Graph DP Problems:}
      \begin{itemize}
        \item Shortest paths (Floyd-Warshall)
        \item Traveling Salesman (TSP)
        \item Longest path in DAG
        \item Number of paths
      \end{itemize}

      \vspace{0.3cm}
      \textbf{DAG Pattern:}
      \begin{enumerate}
        \item Topological sort
        \item Process in topo order
        \item Compute DP for each vertex
      \end{enumerate}
    \end{column}

    \begin{column}{0.5\textwidth}
      \textbf{TSP with Bitmask DP:}
      \begin{itemize}
        \item State: \texttt{dp[mask][i]}
        \item \texttt{mask}: visited vertices
        \item \texttt{i}: current vertex
        \item Time: $O(2^n \cdot n^2)$
      \end{itemize}

      \vspace{0.3cm}
      \textbf{Applications:}
      \begin{itemize}
        \item Route optimization
        \item Circuit board drilling
      \end{itemize}
    \end{column}
  \end{columns}
\end{frame}

%% Complexity-Driven Design
\section{Complexity-Driven Design}

\begin{frame}{Complexity-Driven Design Philosophy}
  \begin{block}{Core Principle}
    Choose data structures and algorithms based on complexity requirements
  \end{block}

  \vspace{0.3cm}
  \textbf{Design Process:}
  \begin{enumerate}
    \item Identify operations needed
    \item Determine frequency of each operation
    \item Analyze required time/space complexity
    \item Select optimal data structure + algorithm
  \end{enumerate}

  \vspace{0.3cm}
  \textbf{Trade-off Considerations:}
  \begin{itemize}
    \item Time vs. space
    \item Preprocessing vs. query time
    \item Average vs. worst-case performance
  \end{itemize}
\end{frame}

\begin{frame}{Complexity Analysis Framework}
  \begin{columns}
    \begin{column}{0.5\textwidth}
      \textbf{Common Complexities:}
      \begin{itemize}
        \item $O(1)$ - Constant
        \item $O(\log n)$ - Logarithmic
        \item $O(n)$ - Linear
        \item $O(n \log n)$ - Linearithmic
        \item $O(n^2)$ - Quadratic
        \item $O(2^n)$ - Exponential
      \end{itemize}
    \end{column}

    \begin{column}{0.5\textwidth}
      \textbf{Growth Comparison:}
      \begin{itemize}
        \item $n = 10$: all fast
        \item $n = 100$: $O(n^2)$ noticeable
        \item $n = 1000$: $O(n \log n)$ max
        \item $n = 10^6$: $O(n)$ or better
        \item $n = 10^9$: $O(\log n)$ or $O(1)$
      \end{itemize}
    \end{column}
  \end{columns}

  \vspace{0.3cm}
  \begin{alertblock}{Critical Insight}
    A faster algorithm with better complexity will eventually outperform a slower one, regardless of constant factors
  \end{alertblock}
\end{frame}

\begin{frame}{Data Structure Selection Guide}
  \begin{table}
    \centering
    \tiny
    \begin{tabular}{lcccc}
      \hline
      \textbf{Need} & \textbf{Insert} & \textbf{Search} & \textbf{Delete} & \textbf{Structure} \\
      \hline
      Fast access & - & $O(1)$ & - & Array / Hash \\
      Fast insert/delete & $O(1)$ & - & $O(1)$ & Linked List \\
      Sorted + search & $O(n)$ & $O(\log n)$ & $O(n)$ & Sorted Array \\
      Sorted + dynamic & $O(\log n)$ & $O(\log n)$ & $O(\log n)$ & BST / Heap \\
      Range queries & - & $O(\log n)$ & - & Segment Tree \\
      Key-value & $O(1)$ & $O(1)$ & $O(1)$ & Hash Table \\
      Priority & $O(\log n)$ & $O(1)$ min & $O(\log n)$ & Heap \\
      \hline
    \end{tabular}
  \end{table}

  \vspace{0.3cm}
  \begin{block}{Selection Criteria}
    \begin{itemize}
      \item Identify the most frequent operation
      \item Optimize for that operation
      \item Accept trade-offs for less frequent operations
    \end{itemize}
  \end{block}
\end{frame}

\begin{frame}{Algorithm Selection Examples}
  \begin{columns}
    \begin{column}{0.5\textwidth}
      \textbf{Sorting Selection:}
      \begin{itemize}
        \item Small array ($n < 50$) $\to$ Insertion
        \item General purpose $\to$ Quick Sort
        \item Stable needed $\to$ Merge Sort
        \item Limited memory $\to$ Heap Sort
        \item Integer range $\to$ Counting Sort
      \end{itemize}
    \end{column}

    \begin{column}{0.5\textwidth}
      \textbf{Graph Algorithm Selection:}
      \begin{itemize}
        \item Single shortest path $\to$ Dijkstra
        \item Negative weights $\to$ Bellman-Ford
        \item All pairs $\to$ Floyd-Warshall
        \item MST sparse $\to$ Kruskal
        \item MST dense $\to$ Prim
      \end{itemize}
    \end{column}
  \end{columns}
\end{frame}

%% Practical Optimization Patterns
\section{Practical Optimization Patterns}

\begin{frame}{Optimization Patterns Overview}
  \begin{block}{What Are Patterns?}
    Reusable techniques that optimize algorithms for specific problem structures
  \end{block}

  \vspace{0.3cm}
  \textbf{Common Patterns:}
  \begin{itemize}
    \item Two Pointers
    \item Sliding Window
    \item Prefix Sum
    \item Monotonic Stack
    \item Binary Search Patterns
    \item Greedy with Sorting
  \end{itemize}

  \vspace{0.3cm}
  \textbf{Benefits:}
  \begin{itemize}
    \item Reduce time complexity (often $O(n^2) \to O(n)$)
    \item Simplify implementation
    \item Widely applicable
  \end{itemize}
\end{frame}

\subsection{Two Pointers}

\begin{frame}{Two Pointers Pattern}
  \begin{columns}
    \begin{column}{0.5\textwidth}
      \textbf{Concept:}
      \begin{itemize}
        \item Use two indices/pointers
        \item Move based on conditions
        \item Avoid nested loops
      \end{itemize}

      \vspace{0.3cm}
      \textbf{Variants:}
      \begin{itemize}
        \item Opposite direction (start/end)
        \item Same direction (slow/fast)
        \item Two arrays
      \end{itemize}
    \end{column}

    \begin{column}{0.5\textwidth}
      \textbf{Common Problems:}
      \begin{itemize}
        \item Two Sum (sorted)
        \item Container with most water
        \item Remove duplicates
        \item Palindrome check
      \end{itemize}

      \vspace{0.3cm}
      \textbf{Complexity:}
      \begin{itemize}
        \item Time: $O(n)$ (single pass)
        \item Space: $O(1)$
      \end{itemize}
    \end{column}
  \end{columns}
\end{frame}

\begin{frame}[fragile]{Two Pointers: Example}
\begin{lstlisting}[language=Python]
def two_sum_sorted(arr, target):
    """Find pair summing to target in sorted array"""
    left, right = 0, len(arr) - 1

    while left < right:
        current_sum = arr[left] + arr[right]

        if current_sum == target:
            return [left, right]
        elif current_sum < target:
            left += 1
        else:
            right -= 1

    return []  # No pair found
\end{lstlisting}

Time: $O(n)$ instead of $O(n^2)$ brute force
\end{frame}

\subsection{Sliding Window}

\begin{frame}{Sliding Window Pattern}
  \begin{columns}
    \begin{column}{0.5\textwidth}
      \textbf{Concept:}
      \begin{itemize}
        \item Maintain a window over data
        \item Expand/contract window
        \item Track window properties
      \end{itemize}

      \vspace{0.3cm}
      \textbf{Types:}
      \begin{itemize}
        \item Fixed size window
        \item Variable size window
      \end{itemize}
    \end{column}

    \begin{column}{0.5\textwidth}
      \textbf{Common Problems:}
      \begin{itemize}
        \item Maximum sum subarray (size $k$)
        \item Longest substring without repeats
        \item Minimum window substring
        \item Subarray product less than $k$
      \end{itemize}

      \vspace{0.3cm}
      \textbf{Complexity:}
      \begin{itemize}
        \item Time: $O(n)$ (each element visited twice max)
        \item Space: $O(k)$ for window state
      \end{itemize}
    \end{column}
  \end{columns}
\end{frame}

\begin{frame}[fragile]{Sliding Window: Fixed Size Example}
\begin{lstlisting}[language=Python]
def max_sum_subarray(arr, k):
    """Find maximum sum of subarray of size k"""
    window_sum = sum(arr[:k])
    max_sum = window_sum

    for i in range(k, len(arr)):
        # Slide window: remove left, add right
        window_sum = window_sum - arr[i - k] + arr[i]
        max_sum = max(max_sum, window_sum)

    return max_sum
\end{lstlisting}

Time: $O(n)$ instead of $O(nk)$ recomputing each window
\end{frame}

\subsection{Prefix Sum}

\begin{frame}{Prefix Sum Pattern}
  \begin{columns}
    \begin{column}{0.5\textwidth}
      \textbf{Concept:}
      \begin{itemize}
        \item Precompute cumulative sums
        \item Answer range queries in $O(1)$
        \item Trade space for time
      \end{itemize}

      \vspace{0.3cm}
      \textbf{Formula:}
      \begin{align*}
        &\text{prefix}[i] = \sum_{j=0}^{i} \text{arr}[j] \\
        &\text{sum}[l, r] = \text{prefix}[r] - \text{prefix}[l-1]
      \end{align*}
    \end{column}

    \begin{column}{0.5\textwidth}
      \textbf{Common Problems:}
      \begin{itemize}
        \item Range sum queries
        \item Subarray sum equals $k$
        \item Equilibrium index
        \item 2D matrix sum queries
      \end{itemize}

      \vspace{0.3cm}
      \textbf{Complexity:}
      \begin{itemize}
        \item Preprocess: $O(n)$
        \item Query: $O(1)$
        \item Space: $O(n)$
      \end{itemize}
    \end{column}
  \end{columns}
\end{frame}

\subsection{Monotonic Stack}

\begin{frame}{Monotonic Stack Pattern}
  \begin{columns}
    \begin{column}{0.5\textwidth}
      \textbf{Concept:}
      \begin{itemize}
        \item Stack maintaining monotonic property
        \item Pop elements violating property
        \item Efficient for next/previous greater/smaller
      \end{itemize}

      \vspace{0.3cm}
      \textbf{Types:}
      \begin{itemize}
        \item Monotonic increasing
        \item Monotonic decreasing
      \end{itemize}
    \end{column}

    \begin{column}{0.5\textwidth}
      \textbf{Common Problems:}
      \begin{itemize}
        \item Next greater element
        \item Largest rectangle in histogram
        \item Stock span problem
        \item Trapping rain water
      \end{itemize}

      \vspace{0.3cm}
      \textbf{Complexity:}
      \begin{itemize}
        \item Time: $O(n)$ (each element pushed/popped once)
        \item Space: $O(n)$
      \end{itemize}
    \end{column}
  \end{columns}
\end{frame}

\subsection{Binary Search Patterns}

\begin{frame}{Binary Search on Answer}
  \begin{columns}
    \begin{column}{0.5\textwidth}
      \textbf{Concept:}
      \begin{itemize}
        \item Search space is answer range
        \item Not searching in array
        \item Check if answer is feasible
      \end{itemize}

      \vspace{0.3cm}
      \textbf{Pattern:}
      \begin{enumerate}
        \item Define search space [low, high]
        \item Binary search on answer
        \item Check feasibility with $O(n)$ function
      \end{enumerate}
    \end{column}

    \begin{column}{0.5\textwidth}
      \textbf{Common Problems:}
      \begin{itemize}
        \item Allocate minimum pages
        \item Split array largest sum
        \item Koko eating bananas
        \item Capacity to ship packages
      \end{itemize}

      \vspace{0.3cm}
      \textbf{Complexity:}
      \begin{itemize}
        \item Time: $O(n \log(\text{range}))$
        \item Space: $O(1)$
      \end{itemize}
    \end{column}
  \end{columns}
\end{frame}

\subsection{Greedy with Sorting}

\begin{frame}{Greedy with Sorting Pattern}
  \begin{columns}
    \begin{column}{0.5\textwidth}
      \textbf{Concept:}
      \begin{itemize}
        \item Sort to reveal greedy structure
        \item Make locally optimal choices
        \item Often requires proof of correctness
      \end{itemize}

      \vspace{0.3cm}
      \textbf{Pattern:}
      \begin{enumerate}
        \item Sort by appropriate criterion
        \item Iterate and make greedy choice
        \item Prove optimal substructure
      \end{enumerate}
    \end{column}

    \begin{column}{0.5\textwidth}
      \textbf{Common Problems:}
      \begin{itemize}
        \item Activity selection
        \item Fractional knapsack
        \item Meeting rooms
        \item Non-overlapping intervals
      \end{itemize}

      \vspace{0.3cm}
      \textbf{Complexity:}
      \begin{itemize}
        \item Time: $O(n \log n)$ for sorting
        \item Space: $O(1)$ or $O(n)$
      \end{itemize}
    \end{column}
  \end{columns}
\end{frame}

\begin{frame}{Optimization Patterns: Summary}
  \begin{table}
    \centering
    \small
    \begin{tabular}{lcc}
      \hline
      \textbf{Pattern} & \textbf{Time Improvement} & \textbf{Common Use} \\
      \hline
      Two Pointers & $O(n^2) \to O(n)$ & Sorted arrays, pairs \\
      Sliding Window & $O(nk) \to O(n)$ & Subarrays, substrings \\
      Prefix Sum & $O(n) \to O(1)$ per query & Range queries \\
      Monotonic Stack & $O(n^2) \to O(n)$ & Next greater/smaller \\
      Binary Search on Answer & $O(n^2) \to O(n \log R)$ & Optimization problems \\
      Greedy + Sorting & Varies & Scheduling, intervals \\
      \hline
    \end{tabular}
  \end{table}

  \vspace{0.3cm}
  \begin{block}{Pattern Recognition}
    Learning these patterns helps identify optimization opportunities in new problems
  \end{block}
\end{frame}

% %% Summary
\section{Summary}

\begin{frame}{Course Summary}
  \textbf{Sorting Algorithms:}
  \begin{itemize}
    \item Comparison-based: $\Omega(n \log n)$ lower bound
    \item Non-comparison: Can beat lower bound with constraints
    \item Choose based on data characteristics and constraints
  \end{itemize}

  \vspace{0.2cm}
  \textbf{Searching Algorithms:}
  \begin{itemize}
    \item Sorted arrays enable $O(\log n)$ search
    \item Hash tables provide $O(1)$ average lookup
    \item Specialized searches for specific scenarios
  \end{itemize}

  \vspace{0.2cm}
  \textbf{Graph Algorithms:}
  \begin{itemize}
    \item Shortest path: Dijkstra, Bellman-Ford, Floyd-Warshall, A*
    \item MST: Kruskal (sparse), Prim (dense)
  \end{itemize}
\end{frame}

\begin{frame}{Key Takeaways}
  \textbf{Dynamic Programming:}
  \begin{itemize}
    \item Data structures enable efficient DP implementation
    \item Arrays for tabulation, hash tables for memoization
    \item Choose structure based on state space
  \end{itemize}

  \vspace{0.2cm}
  \textbf{Complexity-Driven Design:}
  \begin{itemize}
    \item Analyze required operations and frequencies
    \item Select data structures optimizing common operations
    \item Accept trade-offs for rare operations
  \end{itemize}

  \vspace{0.2cm}
  \textbf{Optimization Patterns:}
  \begin{itemize}
    \item Two pointers, sliding window, prefix sum, monotonic stack
    \item Often reduce $O(n^2) \to O(n)$
    \item Pattern recognition is key skill
  \end{itemize}
\end{frame}

\begin{frame}{Final Thoughts}
  \begin{block}{Algorithm + Data Structure = Program}
    Neither algorithms nor data structures exist in isolation. Mastery requires understanding their synergy.
  \end{block}

  \vspace{0.3cm}
  \textbf{Problem-Solving Process:}
  \begin{enumerate}
    \item Understand the problem and constraints
    \item Identify the required complexity
    \item Recognize applicable patterns
    \item Choose appropriate data structure
    \item Implement and optimize
  \end{enumerate}

  \vspace{0.3cm}
  \begin{center}
    \textit{``The best algorithm is the one that solves your specific problem efficiently.''}
  \end{center}
\end{frame}

\begin{frame}
  \begin{center}
    {\Huge Thank You!}

    \vspace{1cm}
    {\Large Questions?}
  \end{center}
\end{frame}

\end{document}
